% III. Проектиране на системата.
% Тук се описва структурата и границите между слоевете, преди какъвто и да е код.
\section{Проектиране на системата}
\label{sec:design}

\subsection{Обща структура}

Системата е разделена на четири слоя. Най-долу стои слой за вход и изход, който
предоставя две операции: изпращане на съобщение до съсед и получаване на съобщение.
Над него стои източник на време, който предоставя текущ момент и планиране на таймер.
Над двата стои протоколното ядро, което съдържа автоматите за състоянията на съседите,
базата с обявленията за състояние на връзките и изчисляването на маршрутите. Най-отгоре
стои приложната обвивка, която в единия режим е демон върху операционната система, а в
другия сценарий на симулатора. Протоколното ядро не съдържа обръщения към часовника и
към сокетите. Това ограничение е причината двата режима да изпълняват един и същ код и
трябва да бъде проверявано, а не да се разчита на дисциплина при писане.
\TODO{да се добави фигура на слоевете с легенда и препратка към нея в текста}

\subsection{Абстракция върху времето и входа и изхода}

В режим на симулация източникът на време е виртуален часовник, който се движи само
когато опашката от събития извади следващото събитие. Между две събития не изтича
време, така че изпълнението не зависи от натовареността на машината. Планирането на
таймер добавя събитие в опашката, а изпращането на съобщение добавя събитие за
получаване след закъснение, определено от модела на връзката. При еднакво начално
състояние и еднаква последователност от външни събития изпълнението дава един и същ
резултат, което е условието за възпроизводимо измерване.

В режим на демон източникът на време е монотонният часовник на операционната система, а
входът и изходът са сокети. Тъй като целевите платформи са Windows, Linux и macOS, слоят
за вход и изход има отделна реализация за всяка от тях зад общ интерфейс, вместо кодът
да се пише срещу интерфейса на една платформа и после да се пренася.
\TODO{да се фиксира механизмът за изчакване на събития по платформи, след измерване дали
разликата между наличните механизми изобщо влияе на резултата}

\subsection{Протоколно ядро}

Ядрото поддържа три структури. Първата е таблица със съседите, всеки от които се намира
в състояние от автомата за съседство и има таймери за поддържане и за отпадане.
Втората е база с обявленията за състояние на връзките, индексирана по подател и с
поредни номера, които определят кое обявление е по-ново. Третата е маршрутната таблица,
получена от изчисляването на дървото от най-кратки пътища върху графа, реконструиран от
базата. Разпространението на обявленията е лавинно: получено по-ново обявление се
записва и се препраща по всички връзки без тази, по която е дошло.

Разделението между база и маршрутна таблица е съществено за измерването. Моментът, в
който възел научава за промяна, и моментът, в който възел започва да препраща по новия
път, са различни, а разликата между тях е точно интервалът, в който възникват преходни
цикли. Ако двете структури бяха обединени, тази разлика би била ненаблюдаема.

\subsection{Записване на събития за анализ}

Всяко значимо събитие, тоест промяна на състоянието на съсед, приемане на обявление,
изчисляване на маршрути и инсталиране на маршрут, се записва с виртуалния или реалния
момент на настъпването си в поток от събития. Анализът след изпълнението работи само
върху този поток и не се нуждае от достъп до вътрешното състояние. Така една и съща
процедура за анализ обработва изхода и от двата режима.
\TODO{да се фиксира форматът на записа на събитията и да се приложи като приложение}
